\section{结语：北方危机是关系的重组}

北宋的终结不能只归结为开封的失守。它的中枢、漕运、市场与文书曾使辽、西夏和北方边地的压力能够被纳入预算和行政，却也使燕云、河北与西北的军事竞争持续牵动东南财赋和朝廷政治。宋与辽的和议、西夏边境的战守，以及后来与金的联盟和冲突，都是各方以不同资源组合处理风险的选择，而非北宋单方面的外交成败。

靖康使这种组合断裂，却没有清空其社会基础。北迁者、南渡者、接收华北的金朝官民和留在城市、乡村与水道上的居民，分别携带并改造了旧有的制度与关系。下一章因此不把辽、西夏写作宋亡的背景，而把它们作为拥有自身财政、军政、宗教和边境选择的政治体来考察。
